% One-pager digest — Milano complex-frequency PSCC 2026 tutorial.
\documentclass[10pt]{article}
\usepackage[a4paper,margin=1.4cm]{geometry}
\usepackage{amsmath,amssymb,bm}
\usepackage{multicol}
\usepackage{enumitem}
\setlist{nosep,leftmargin=1.2em}
\usepackage{titlesec}
\titleformat{\section}{\bfseries\small}{}{0pt}{}
\titlespacing{\section}{0pt}{4pt}{2pt}
\pagestyle{empty}
\setlength{\parindent}{0pt}

\newcommand{\cb}[1]{\bar{#1}}

\begin{document}
{\large\bfseries Complex Frequency, Local Synchronization \& Transient Slack Capability}\\[-2pt]
{\footnotesize PSCC 2026 Tutorial \hfill Federico Milano (UC Dublin) \hfill Low-inertia power systems}
\hrule\vspace{4pt}

\section*{In one line}
One complex quantity $\cb\eta=\rho+j\omega$ fuses the rate of change of a bus voltage's
\emph{magnitude} ($\rho$) and its \emph{angular frequency} ($\omega$), giving a single
language for the dynamics, stability and synchronization of converter-dominated grids.

\begin{multicols}{2}
\section*{Key ideas}
\begin{itemize}
\item \textbf{Complex frequency:} with $\cb v = v\,e^{\,u+j\theta}$ ($u=\ln v$),
  $$\cb\eta=\tfrac{d}{dt}(u+j\theta)=\rho+j\omega,\quad \rho=\dot v/v,\ \omega=\dot\theta.$$
\item \textbf{Complex power} $\cb s=\cb v\circ\cb\imath^{*}=p+jq$, with rate of change (RoCoP)
  $\cb s' = \cb\eta\,\cb s + \cb v\circ\cb\imath'^{*}$.
\item \textbf{Local synchronization:} a grid-side view of synchronism --- device is
  synchronized when $\rho\!\to\!0$ and $\omega\!\to\!\omega_r$ (bounded/asymptotic, BLS/ALS).
\item \textbf{Transient Slack Capability (TSC):} what a device needs to sustain the grid
  after a large disturbance; realised by the \textbf{Dual grid-forming converter}.
\end{itemize}

\section*{Concept map}
Park vectors $\to$ $\cb\eta$ (complex freq.) $\to$ RoCoP $\cb s'$ $\to$ device component
$\cb\chi=\cb\xi-\cb\eta$ $\to$ local sync (BLS/ALS) $\to$ TSC $\to$ Dual-GFM application.

\section*{Worked takeaways (verified)}
\begin{itemize}
\item $\cb v' = \cb v\circ\cb\eta$;\ \ $\cb\imath'=\cb\eta\,\cb\imath$ (constant $\cb Y$).
\item RoCoP: $\cb s' = \cb\eta\,\cb s + \cb v\circ\cb\imath'^{*}$.
\item General component: $\cb\imath'=\cb\imath(\cb\chi+\cb\eta)\Rightarrow \cb\chi=\cb\xi-\cb\eta$.
\end{itemize}
\emph{Why $\rho,\omega$ combine:} both are $[\text{time}]^{-1}$ (see \texttt{verify.md}).

\section*{Prerequisites (your profile)}
Park/dq transform \& Hadamard $\circ$ (\texttt{seen}); complex power (\texttt{followed});
target for $\cb\eta$ itself: \texttt{followed}$\to$\texttt{intuitive}.

\section*{Where it fits}
Milano's programme: \emph{Complex Frequency} (TPWRS 2022, arXiv:2105.07769) $\to$
\emph{Local Synchronization} (Ponce \& Milano 2025) $\to$ \emph{Dual-GFM} \&
\emph{Transient Slack Capability} (2025).

\section*{Relevance to my work}
Speaks to \texttt{effective\_inertia}: $\cb\eta=\rho+j\omega$ and RoCoP split power-rate into
voltage ($\rho$) and frequency ($\omega$) parts --- a language for the fast-response
\emph{offset} beyond $H_{\rm load}$ and for $r=H_{\rm eff}/H_{\rm load}$ (Milano's
stored-vs-delivered gap). $\rho=\dot v/v$ is exactly your CMLD voltage-response / feeder
question; $\cb\chi=\cb\xi-\cb\eta$ + local sync formalise the motor slip--torque coupling.
Continuous framing --- \emph{not} \texttt{pv\_trip}'s discrete 49.5/49\,Hz switching.
\end{multicols}
\end{document}
